\documentclass[11pt]{article}

\usepackage[margin=0.85in]{geometry}
\usepackage[T1]{fontenc}
\usepackage[utf8]{inputenc}
\usepackage{amsmath,amssymb}
\usepackage{xcolor}
\usepackage{booktabs}
\usepackage{tabularx}
\usepackage{tcolorbox}

\definecolor{auditblue}{RGB}{34,80,145}
\definecolor{auditred}{RGB}{185,38,38}
\definecolor{auditgray}{RGB}{245,246,248}

\setlength{\parindent}{0pt}
\setlength{\parskip}{0.55em}

\newtcolorbox{resultbox}[1]{
  colback=auditgray,
  colframe=#1,
  boxrule=0.8pt,
  arc=1.5pt,
  left=7pt,
  right=7pt,
  top=6pt,
  bottom=6pt
}

\begin{document}

\begin{center}
{\Large\bfseries Real-\(z\) Check of the Shortcut Sign in Eq. (29)}\\[0.3em]
{\small Comparison with the direct stochastic shortcut matrix}
\end{center}

\begin{resultbox}{auditblue}
\textbf{Short answer.}
The quick check locates the sign issue already in Eq. (29).  Used literally,
the printed Eq. (29) reproduces the original Eq. (32), but both disagree with
the direct stochastic shortcut matrix.  With the stochastic shortcut sign, the
Eq. (29) ratio and the corrected Eq. (32) agree with the direct matrix to
roundoff.
\end{resultbox}

\section*{Direct Matrix Benchmark}

The benchmark is the finite-state transition matrix for the stochastic shortcut
model itself, not Eq. (29) or Eq. (32).  Start from the lazy ring matrix
\[
T_0(i,i)=1-q,\qquad T_0(i,i\pm1)=q/2,
\]
and move the shortcut mass
\[
\lambda=\beta(1-q)
\]
from the \(u\to u\) self-loop to the directed edge \(u\to v\):
\[
\widetilde T=T_0+\lambda(E_{uv}-E_{uu}).
\]
Then compute the resolvent and the first-passage ratio directly:
\[
\widetilde S(z)=(I-z\widetilde T)^{-1},\qquad
\widetilde F_{n_0}(v,z)=
\frac{\widetilde S_{n_0v}(z)}{\widetilde S_{vv}(z)}.
\]

\section*{Sign Implied by the Matrix}

The direct stochastic matrix gives
\[
I-z\widetilde T=(I-zT_0)+z\lambda(E_{uu}-E_{uv}).
\]
Writing \(P(z)=(I-zT_0)^{-1}\), the corresponding rank-one update has the
sign pattern
\[
\widetilde P_{ij}(z)
=P_{ij}(z)
-\frac{z\lambda P_{iu}(z)\{P_{uj}(z)-P_{vj}(z)\}}
{1+z\lambda\{P_{uu}(z)-P_{vu}(z)\}}.
\]
Thus the stochastic shortcut \(u\to v\) gives a plus sign in this denominator
and the corresponding minus sign in the correction term.  The printed Eq. (29)
has the opposite pattern: a plus sign in the correction term and a
\(1-z\lambda\{P_{uu}(z)-P_{vu}(z)\}\) denominator.  That is the inverse for the
opposite rank-one defect, not for removing probability mass from \(u\to u\) and
adding it to \(u\to v\).  This is why the sign issue is already in Eq. (29).

\section*{Numerical Check}

The grid uses the figure parameters \(N=10\), \(q=2/3\), \(\beta=2/7\),
\(u=5\), \(v=0\), plus a second symmetric case, with several starts \(n_0\)
and real values of \(z\).  Site labels here are zero-based.

\subsection*{Main sign check}

\begin{center}
\begin{tabularx}{0.98\linewidth}{@{}Xr@{}}
\toprule
Quantity compared over the grid & Maximum absolute error \\
\midrule
Literal printed Eq. (29) ratio vs direct stochastic matrix & \(3.669746\times10^{-1}\) \\
Original Eq. (32) vs direct stochastic matrix & \(3.669746\times10^{-1}\) \\
Stochastic-sign Eq. (29) ratio vs direct stochastic matrix & \(2.220446\times10^{-16}\) \\
Corrected Eq. (32) vs direct stochastic matrix & \(2.775558\times10^{-16}\) \\
Literal printed Eq. (29) ratio vs original Eq. (32) & \(1.318390\times10^{-16}\) \\
\bottomrule
\end{tabularx}
\end{center}

\subsection*{Control checks}

\begin{center}
\begin{tabularx}{0.98\linewidth}{@{}Xr@{}}
\toprule
Quantity compared over the grid & Maximum absolute error \\
\midrule
Eq. (19) \(W_{n_0}(u,z)\) vs absorbing matrix & \(1.043610\times10^{-14}\) \\
Eq. (28) \(W_u(u,z)\) vs absorbing matrix & \(1.687539\times10^{-14}\) \\
\bottomrule
\end{tabularx}
\end{center}

The control checks show that the discrepancy is not coming from the
\(W_{n_0}(u,z)\) or \(W_u(u,z)\) inputs used in Eq. (32).

\section*{Representative Values}

\begin{center}
\small
\begin{tabular}{@{}rrrrr@{}}
\toprule
\(n_0\) & \(z\) & direct matrix & printed Eq. (29) / original Eq. (32)
& stochastic sign / corrected Eq. (32) \\
\midrule
2 & 0.80 & 0.197025316799 & 0.173752283343 & 0.197025316799 \\
4 & 0.80 & 0.095953324694 & -0.033502923905 & 0.095953324694 \\
\bottomrule
\end{tabular}
\end{center}

\begin{resultbox}{auditred}
\textbf{Interpretation.}
For \(0<z<1\), a first-passage generating function has non-negative
coefficients, so it should not become negative.  The negative value in the
second row is a useful sanity check that the printed sign is not the stochastic
shortcut sign for this \(u\to v\) model.
\end{resultbox}

\section*{Conclusion}

The printed Eq. (29) and the original Eq. (32) are internally consistent with
each other.  The problem is that this shared sign convention differs from the
direct probability-conserving stochastic shortcut matrix.  Therefore the sign
issue is already present in Eq. (29) under the stochastic \(u\to v\)
interpretation.

\end{document}
